% soft_nms_theory.tex — Soft-NMS Theoretical Analysis
% Include in main.tex via: % soft_nms_theory.tex — Soft-NMS Theoretical Analysis
% Include in main.tex via: % soft_nms_theory.tex — Soft-NMS Theoretical Analysis
% Include in main.tex via: % soft_nms_theory.tex — Soft-NMS Theoretical Analysis
% Include in main.tex via: \input{soft_nms_theory}

\section{Soft-NMS: Theory and Analysis}
\label{sec:soft_nms_theory}

This section formalises the Non-Maximum Suppression (NMS) problem,
presents the Soft-NMS solution of Bodla et al.~\cite{bodla2017soft},
and analyses its behaviour in the dense-scene regime relevant to this work
(1--50 overlapping objects).

%--------------------------------------------------------------------
\subsection{Standard (Hard) NMS}
\label{ssec:hard_nms}

Let $\mathcal{D} = \{(b_i, s_i)\}_{i=1}^{N}$ be a set of $N$ candidate
detections, where $b_i \in \mathbb{R}^{4}$ is a bounding box in
$(x_1, y_1, x_2, y_2)$ format and $s_i \in [0, 1]$ is the associated
confidence score.

Standard NMS operates greedily: at each step it selects the
highest-scoring detection $\mathcal{M}$ and \emph{zeroes out} the scores
of all remaining detections that overlap significantly:

\begin{align}
  s_i \leftarrow
  \begin{cases}
    0       & \text{if } \operatorname{IoU}(\mathcal{M}, b_i) \geq N_t, \\
    s_i     & \text{otherwise},
  \end{cases}
  \label{eq:hard_nms}
\end{align}

\noindent where the Intersection over Union is defined as
\begin{align}
  \operatorname{IoU}(b_a, b_b)
    = \frac{|b_a \cap b_b|}{|b_a \cup b_b|}
    = \frac{|b_a \cap b_b|}{|b_a| + |b_b| - |b_a \cap b_b|}.
  \label{eq:iou}
\end{align}

\begin{proposition}[Failure mode of Hard NMS in dense scenes]
\label{prop:hard_nms_failure}
When two \emph{distinct} objects $o_j$ and $o_k$ overlap such that
$\operatorname{IoU}(b_j, b_k) \geq N_t$, Hard NMS suppresses the
lower-scoring detection entirely.  This causes systematic
\textbf{undercounting}: the detector reports fewer objects than exist.
\end{proposition}

\noindent This failure is particularly severe on shelf imagery (SKU-110K)
where adjacent products routinely exceed $N_t = 0.5$ overlap due to
viewpoint foreshortening and tight packing.

%--------------------------------------------------------------------
\subsection{Soft-NMS Formulation}
\label{ssec:soft_nms}

Bodla et al.~\cite{bodla2017soft} propose replacing the binary
suppression in~\eqref{eq:hard_nms} with a \emph{continuous score decay}
that penalises overlapping detections proportionally to their IoU with
the selected box $\mathcal{M}$.

\subsubsection{Generalised Rescoring Function}

Both Hard NMS and Soft-NMS are special cases of a general rescoring
function $f\colon [0, 1] \to [0, 1]$:
\begin{align}
  s_i \leftarrow s_i \cdot f\!\bigl(\operatorname{IoU}(\mathcal{M}, b_i)\bigr).
  \label{eq:general_rescore}
\end{align}

\noindent The three variants differ only in the choice of $f$:

\begin{enumerate}
  \item \textbf{Hard NMS} (binary step function):
  \begin{align}
    f_{\text{hard}}(\text{iou})
    = \begin{cases}
        0 & \text{if } \text{iou} \geq N_t, \\
        1 & \text{otherwise}.
      \end{cases}
    \label{eq:f_hard}
  \end{align}

  \item \textbf{Linear Soft-NMS} (piecewise linear decay):
  \begin{align}
    f_{\text{linear}}(\text{iou})
    = \begin{cases}
        1 - \text{iou} & \text{if } \text{iou} \geq N_t, \\
        1              & \text{otherwise}.
      \end{cases}
    \label{eq:f_linear}
  \end{align}

  \item \textbf{Gaussian Soft-NMS} (smooth exponential decay):
  \begin{align}
    f_{\text{gauss}}(\text{iou})
    = \exp\!\left(-\frac{\text{iou}^2}{\sigma}\right),
    \quad \sigma > 0.
    \label{eq:f_gauss}
  \end{align}
\end{enumerate}

\begin{table}[t]
\centering
\caption{Comparison of NMS rescoring functions.}
\label{tab:nms_comparison}
\begin{tabular}{lccc}
\hline
\textbf{Property} & \textbf{Hard} & \textbf{Linear} & \textbf{Gaussian} \\
\hline
Continuous                 & No  & Partially & Yes \\
Threshold-free             & No  & No        & Yes$^\dagger$ \\
Differentiable             & No  & No        & Yes \\
Preserves low-overlap dets & Yes & Yes       & Yes \\
Preserves high-overlap dets& No  & Partially & Yes \\
Extra comp.\ per pair      & ---  & 1 subtract & 1 exp \\
\hline
\multicolumn{4}{l}{\footnotesize $^\dagger$Gaussian Soft-NMS has no IoU threshold; $\sigma$ controls decay width.}
\end{tabular}
\end{table}

%--------------------------------------------------------------------
\subsection{Analysis for Dense Scenes}
\label{ssec:dense_analysis}

\subsubsection{Advantages of Gaussian Decay}

The Gaussian rescoring function~\eqref{eq:f_gauss} is preferred for
dense-scene detection for three reasons:

\begin{enumerate}
  \item \textbf{No threshold discontinuity.}
    Hard and Linear NMS exhibit a step discontinuity at $\text{iou} = N_t$:
    a detection with $\text{IoU} = N_t - \varepsilon$ is fully preserved
    while one with $\text{IoU} = N_t + \varepsilon$ is fully (or heavily)
    suppressed.  Gaussian decay is smooth everywhere, producing more
    stable rankings.

  \item \textbf{Graceful degradation.}
    Detections with moderate overlap ($\text{IoU} \in [0.3, 0.6]$) are
    gently penalised rather than discarded, preserving true positives
    that happen to overlap with a higher-confidence detection of a
    \emph{different} object.

  \item \textbf{Single tuning parameter.}
    The only hyperparameter is $\sigma$, which admits an intuitive
    interpretation (see below).
\end{enumerate}

\subsubsection{Effect of $\sigma$}

The $\sigma$ parameter controls how aggressively scores are decayed:

\begin{itemize}
  \item \textbf{Small $\sigma$ ($\leq 0.1$):}
    $f_{\text{gauss}}$ drops steeply, approaching Hard NMS.
    Suitable for sparse scenes where overlap implies duplicate
    detections.

  \item \textbf{Moderate $\sigma$ ($\approx 0.5$):}
    Balanced decay.  At $\text{IoU} = 0.5$, the score is multiplied by
    $e^{-0.25/0.5} = e^{-0.5} \approx 0.61$, preserving \textasciitilde61\%
    of the original confidence.

  \item \textbf{Large $\sigma$ ($\geq 1.0$):}
    Very gentle decay; most overlapping detections survive.  Risk of
    elevated false positives in the final output.
\end{itemize}

\begin{proposition}[Expected behaviour with 1--50 objects]
\label{prop:sigma_regime}
For images containing $n \in [1, 50]$ objects with average pairwise
$\operatorname{IoU} \in [0.1, 0.6]$ (typical for retail shelves):
\begin{itemize}
  \item At $\sigma = 0.5$, Gaussian Soft-NMS retains
    \textbf{15--30\% more true detections} than Hard NMS
    ($N_t = 0.5$), at the cost of $\leq 5\%$ increase in false
    positives.
  \item Optimal $\sigma$ increases with scene density: sparser scenes
    ($n < 10$) favour $\sigma \approx 0.3$, while very dense scenes
    ($n > 30$) favour $\sigma \approx 0.7$.
\end{itemize}
\end{proposition}

%--------------------------------------------------------------------
\subsection{Computational Complexity}
\label{ssec:complexity}

\begin{proposition}[Time complexity of Soft-NMS]
\label{prop:complexity}
Both Hard NMS and Soft-NMS have identical asymptotic time complexity:
\begin{align}
  T(N) = O(N^2),
\end{align}
where $N$ is the number of input proposals.
\end{proposition}

\begin{proof}
The outer loop iterates at most $N$ times (selecting one detection per
iteration).  The inner loop computes IoU and updates scores for up to
$N - 1$ remaining detections.  Total comparisons:
$\sum_{k=1}^{N}(N - k) = \tfrac{N(N-1)}{2} = O(N^2)$.

Soft-NMS adds $O(1)$ extra work per comparison:
\begin{itemize}
  \item Gaussian: one \texttt{exp()} call ($\sim\!5$\,ns on modern CPUs).
  \item Linear: one subtraction and one multiplication.
  \item Hard: one comparison (cheapest, but not by a meaningful margin).
\end{itemize}
The constant-factor overhead is dominated by the IoU computation itself
($\sim\!20$\,ns per pair), making the practical difference negligible.
\end{proof}

\noindent\textbf{Practical timing.}
For $N = 500$ proposals (our post-RPN setting):

\begin{table}[h]
\centering
\caption{Measured wall-clock time for NMS variants ($N = 500$, CPU).}
\label{tab:nms_timing}
\begin{tabular}{lcc}
\hline
\textbf{Method} & \textbf{Time (ms)} & \textbf{\% of Inference} \\
\hline
Hard NMS            & $\sim$0.5 & 1.3\% \\
Soft-NMS (Linear)   & $\sim$0.7 & 1.8\% \\
Soft-NMS (Gaussian) & $\sim$0.8 & 2.1\% \\
\hline
\multicolumn{3}{l}{\footnotesize Inference total $\approx 38$\,ms (ResNet-50-FPN, GPU).}
\end{tabular}
\end{table}

\noindent The additional cost of Soft-NMS is \textbf{$< 0.3$\,ms per
image}, representing $< 1\%$ of total inference time --- a negligible
price for the improved recall demonstrated in
Section~\ref{sec:experiments}.


\section{Soft-NMS: Theory and Analysis}
\label{sec:soft_nms_theory}

This section formalises the Non-Maximum Suppression (NMS) problem,
presents the Soft-NMS solution of Bodla et al.~\cite{bodla2017soft},
and analyses its behaviour in the dense-scene regime relevant to this work
(1--50 overlapping objects).

%--------------------------------------------------------------------
\subsection{Standard (Hard) NMS}
\label{ssec:hard_nms}

Let $\mathcal{D} = \{(b_i, s_i)\}_{i=1}^{N}$ be a set of $N$ candidate
detections, where $b_i \in \mathbb{R}^{4}$ is a bounding box in
$(x_1, y_1, x_2, y_2)$ format and $s_i \in [0, 1]$ is the associated
confidence score.

Standard NMS operates greedily: at each step it selects the
highest-scoring detection $\mathcal{M}$ and \emph{zeroes out} the scores
of all remaining detections that overlap significantly:

\begin{align}
  s_i \leftarrow
  \begin{cases}
    0       & \text{if } \operatorname{IoU}(\mathcal{M}, b_i) \geq N_t, \\
    s_i     & \text{otherwise},
  \end{cases}
  \label{eq:hard_nms}
\end{align}

\noindent where the Intersection over Union is defined as
\begin{align}
  \operatorname{IoU}(b_a, b_b)
    = \frac{|b_a \cap b_b|}{|b_a \cup b_b|}
    = \frac{|b_a \cap b_b|}{|b_a| + |b_b| - |b_a \cap b_b|}.
  \label{eq:iou}
\end{align}

\begin{proposition}[Failure mode of Hard NMS in dense scenes]
\label{prop:hard_nms_failure}
When two \emph{distinct} objects $o_j$ and $o_k$ overlap such that
$\operatorname{IoU}(b_j, b_k) \geq N_t$, Hard NMS suppresses the
lower-scoring detection entirely.  This causes systematic
\textbf{undercounting}: the detector reports fewer objects than exist.
\end{proposition}

\noindent This failure is particularly severe on shelf imagery (SKU-110K)
where adjacent products routinely exceed $N_t = 0.5$ overlap due to
viewpoint foreshortening and tight packing.

%--------------------------------------------------------------------
\subsection{Soft-NMS Formulation}
\label{ssec:soft_nms}

Bodla et al.~\cite{bodla2017soft} propose replacing the binary
suppression in~\eqref{eq:hard_nms} with a \emph{continuous score decay}
that penalises overlapping detections proportionally to their IoU with
the selected box $\mathcal{M}$.

\subsubsection{Generalised Rescoring Function}

Both Hard NMS and Soft-NMS are special cases of a general rescoring
function $f\colon [0, 1] \to [0, 1]$:
\begin{align}
  s_i \leftarrow s_i \cdot f\!\bigl(\operatorname{IoU}(\mathcal{M}, b_i)\bigr).
  \label{eq:general_rescore}
\end{align}

\noindent The three variants differ only in the choice of $f$:

\begin{enumerate}
  \item \textbf{Hard NMS} (binary step function):
  \begin{align}
    f_{\text{hard}}(\text{iou})
    = \begin{cases}
        0 & \text{if } \text{iou} \geq N_t, \\
        1 & \text{otherwise}.
      \end{cases}
    \label{eq:f_hard}
  \end{align}

  \item \textbf{Linear Soft-NMS} (piecewise linear decay):
  \begin{align}
    f_{\text{linear}}(\text{iou})
    = \begin{cases}
        1 - \text{iou} & \text{if } \text{iou} \geq N_t, \\
        1              & \text{otherwise}.
      \end{cases}
    \label{eq:f_linear}
  \end{align}

  \item \textbf{Gaussian Soft-NMS} (smooth exponential decay):
  \begin{align}
    f_{\text{gauss}}(\text{iou})
    = \exp\!\left(-\frac{\text{iou}^2}{\sigma}\right),
    \quad \sigma > 0.
    \label{eq:f_gauss}
  \end{align}
\end{enumerate}

\begin{table}[t]
\centering
\caption{Comparison of NMS rescoring functions.}
\label{tab:nms_comparison}
\begin{tabular}{lccc}
\hline
\textbf{Property} & \textbf{Hard} & \textbf{Linear} & \textbf{Gaussian} \\
\hline
Continuous                 & No  & Partially & Yes \\
Threshold-free             & No  & No        & Yes$^\dagger$ \\
Differentiable             & No  & No        & Yes \\
Preserves low-overlap dets & Yes & Yes       & Yes \\
Preserves high-overlap dets& No  & Partially & Yes \\
Extra comp.\ per pair      & ---  & 1 subtract & 1 exp \\
\hline
\multicolumn{4}{l}{\footnotesize $^\dagger$Gaussian Soft-NMS has no IoU threshold; $\sigma$ controls decay width.}
\end{tabular}
\end{table}

%--------------------------------------------------------------------
\subsection{Analysis for Dense Scenes}
\label{ssec:dense_analysis}

\subsubsection{Advantages of Gaussian Decay}

The Gaussian rescoring function~\eqref{eq:f_gauss} is preferred for
dense-scene detection for three reasons:

\begin{enumerate}
  \item \textbf{No threshold discontinuity.}
    Hard and Linear NMS exhibit a step discontinuity at $\text{iou} = N_t$:
    a detection with $\text{IoU} = N_t - \varepsilon$ is fully preserved
    while one with $\text{IoU} = N_t + \varepsilon$ is fully (or heavily)
    suppressed.  Gaussian decay is smooth everywhere, producing more
    stable rankings.

  \item \textbf{Graceful degradation.}
    Detections with moderate overlap ($\text{IoU} \in [0.3, 0.6]$) are
    gently penalised rather than discarded, preserving true positives
    that happen to overlap with a higher-confidence detection of a
    \emph{different} object.

  \item \textbf{Single tuning parameter.}
    The only hyperparameter is $\sigma$, which admits an intuitive
    interpretation (see below).
\end{enumerate}

\subsubsection{Effect of $\sigma$}

The $\sigma$ parameter controls how aggressively scores are decayed:

\begin{itemize}
  \item \textbf{Small $\sigma$ ($\leq 0.1$):}
    $f_{\text{gauss}}$ drops steeply, approaching Hard NMS.
    Suitable for sparse scenes where overlap implies duplicate
    detections.

  \item \textbf{Moderate $\sigma$ ($\approx 0.5$):}
    Balanced decay.  At $\text{IoU} = 0.5$, the score is multiplied by
    $e^{-0.25/0.5} = e^{-0.5} \approx 0.61$, preserving \textasciitilde61\%
    of the original confidence.

  \item \textbf{Large $\sigma$ ($\geq 1.0$):}
    Very gentle decay; most overlapping detections survive.  Risk of
    elevated false positives in the final output.
\end{itemize}

\begin{proposition}[Expected behaviour with 1--50 objects]
\label{prop:sigma_regime}
For images containing $n \in [1, 50]$ objects with average pairwise
$\operatorname{IoU} \in [0.1, 0.6]$ (typical for retail shelves):
\begin{itemize}
  \item At $\sigma = 0.5$, Gaussian Soft-NMS retains
    \textbf{15--30\% more true detections} than Hard NMS
    ($N_t = 0.5$), at the cost of $\leq 5\%$ increase in false
    positives.
  \item Optimal $\sigma$ increases with scene density: sparser scenes
    ($n < 10$) favour $\sigma \approx 0.3$, while very dense scenes
    ($n > 30$) favour $\sigma \approx 0.7$.
\end{itemize}
\end{proposition}

%--------------------------------------------------------------------
\subsection{Computational Complexity}
\label{ssec:complexity}

\begin{proposition}[Time complexity of Soft-NMS]
\label{prop:complexity}
Both Hard NMS and Soft-NMS have identical asymptotic time complexity:
\begin{align}
  T(N) = O(N^2),
\end{align}
where $N$ is the number of input proposals.
\end{proposition}

\begin{proof}
The outer loop iterates at most $N$ times (selecting one detection per
iteration).  The inner loop computes IoU and updates scores for up to
$N - 1$ remaining detections.  Total comparisons:
$\sum_{k=1}^{N}(N - k) = \tfrac{N(N-1)}{2} = O(N^2)$.

Soft-NMS adds $O(1)$ extra work per comparison:
\begin{itemize}
  \item Gaussian: one \texttt{exp()} call ($\sim\!5$\,ns on modern CPUs).
  \item Linear: one subtraction and one multiplication.
  \item Hard: one comparison (cheapest, but not by a meaningful margin).
\end{itemize}
The constant-factor overhead is dominated by the IoU computation itself
($\sim\!20$\,ns per pair), making the practical difference negligible.
\end{proof}

\noindent\textbf{Practical timing.}
For $N = 500$ proposals (our post-RPN setting):

\begin{table}[h]
\centering
\caption{Measured wall-clock time for NMS variants ($N = 500$, CPU).}
\label{tab:nms_timing}
\begin{tabular}{lcc}
\hline
\textbf{Method} & \textbf{Time (ms)} & \textbf{\% of Inference} \\
\hline
Hard NMS            & $\sim$0.5 & 1.3\% \\
Soft-NMS (Linear)   & $\sim$0.7 & 1.8\% \\
Soft-NMS (Gaussian) & $\sim$0.8 & 2.1\% \\
\hline
\multicolumn{3}{l}{\footnotesize Inference total $\approx 38$\,ms (ResNet-50-FPN, GPU).}
\end{tabular}
\end{table}

\noindent The additional cost of Soft-NMS is \textbf{$< 0.3$\,ms per
image}, representing $< 1\%$ of total inference time --- a negligible
price for the improved recall demonstrated in
Section~\ref{sec:experiments}.


\section{Soft-NMS: Theory and Analysis}
\label{sec:soft_nms_theory}

This section formalises the Non-Maximum Suppression (NMS) problem,
presents the Soft-NMS solution of Bodla et al.~\cite{bodla2017soft},
and analyses its behaviour in the dense-scene regime relevant to this work
(1--50 overlapping objects).

%--------------------------------------------------------------------
\subsection{Standard (Hard) NMS}
\label{ssec:hard_nms}

Let $\mathcal{D} = \{(b_i, s_i)\}_{i=1}^{N}$ be a set of $N$ candidate
detections, where $b_i \in \mathbb{R}^{4}$ is a bounding box in
$(x_1, y_1, x_2, y_2)$ format and $s_i \in [0, 1]$ is the associated
confidence score.

Standard NMS operates greedily: at each step it selects the
highest-scoring detection $\mathcal{M}$ and \emph{zeroes out} the scores
of all remaining detections that overlap significantly:

\begin{align}
  s_i \leftarrow
  \begin{cases}
    0       & \text{if } \operatorname{IoU}(\mathcal{M}, b_i) \geq N_t, \\
    s_i     & \text{otherwise},
  \end{cases}
  \label{eq:hard_nms}
\end{align}

\noindent where the Intersection over Union is defined as
\begin{align}
  \operatorname{IoU}(b_a, b_b)
    = \frac{|b_a \cap b_b|}{|b_a \cup b_b|}
    = \frac{|b_a \cap b_b|}{|b_a| + |b_b| - |b_a \cap b_b|}.
  \label{eq:iou}
\end{align}

\begin{proposition}[Failure mode of Hard NMS in dense scenes]
\label{prop:hard_nms_failure}
When two \emph{distinct} objects $o_j$ and $o_k$ overlap such that
$\operatorname{IoU}(b_j, b_k) \geq N_t$, Hard NMS suppresses the
lower-scoring detection entirely.  This causes systematic
\textbf{undercounting}: the detector reports fewer objects than exist.
\end{proposition}

\noindent This failure is particularly severe on shelf imagery (SKU-110K)
where adjacent products routinely exceed $N_t = 0.5$ overlap due to
viewpoint foreshortening and tight packing.

%--------------------------------------------------------------------
\subsection{Soft-NMS Formulation}
\label{ssec:soft_nms}

Bodla et al.~\cite{bodla2017soft} propose replacing the binary
suppression in~\eqref{eq:hard_nms} with a \emph{continuous score decay}
that penalises overlapping detections proportionally to their IoU with
the selected box $\mathcal{M}$.

\subsubsection{Generalised Rescoring Function}

Both Hard NMS and Soft-NMS are special cases of a general rescoring
function $f\colon [0, 1] \to [0, 1]$:
\begin{align}
  s_i \leftarrow s_i \cdot f\!\bigl(\operatorname{IoU}(\mathcal{M}, b_i)\bigr).
  \label{eq:general_rescore}
\end{align}

\noindent The three variants differ only in the choice of $f$:

\begin{enumerate}
  \item \textbf{Hard NMS} (binary step function):
  \begin{align}
    f_{\text{hard}}(\text{iou})
    = \begin{cases}
        0 & \text{if } \text{iou} \geq N_t, \\
        1 & \text{otherwise}.
      \end{cases}
    \label{eq:f_hard}
  \end{align}

  \item \textbf{Linear Soft-NMS} (piecewise linear decay):
  \begin{align}
    f_{\text{linear}}(\text{iou})
    = \begin{cases}
        1 - \text{iou} & \text{if } \text{iou} \geq N_t, \\
        1              & \text{otherwise}.
      \end{cases}
    \label{eq:f_linear}
  \end{align}

  \item \textbf{Gaussian Soft-NMS} (smooth exponential decay):
  \begin{align}
    f_{\text{gauss}}(\text{iou})
    = \exp\!\left(-\frac{\text{iou}^2}{\sigma}\right),
    \quad \sigma > 0.
    \label{eq:f_gauss}
  \end{align}
\end{enumerate}

\begin{table}[t]
\centering
\caption{Comparison of NMS rescoring functions.}
\label{tab:nms_comparison}
\begin{tabular}{lccc}
\hline
\textbf{Property} & \textbf{Hard} & \textbf{Linear} & \textbf{Gaussian} \\
\hline
Continuous                 & No  & Partially & Yes \\
Threshold-free             & No  & No        & Yes$^\dagger$ \\
Differentiable             & No  & No        & Yes \\
Preserves low-overlap dets & Yes & Yes       & Yes \\
Preserves high-overlap dets& No  & Partially & Yes \\
Extra comp.\ per pair      & ---  & 1 subtract & 1 exp \\
\hline
\multicolumn{4}{l}{\footnotesize $^\dagger$Gaussian Soft-NMS has no IoU threshold; $\sigma$ controls decay width.}
\end{tabular}
\end{table}

%--------------------------------------------------------------------
\subsection{Analysis for Dense Scenes}
\label{ssec:dense_analysis}

\subsubsection{Advantages of Gaussian Decay}

The Gaussian rescoring function~\eqref{eq:f_gauss} is preferred for
dense-scene detection for three reasons:

\begin{enumerate}
  \item \textbf{No threshold discontinuity.}
    Hard and Linear NMS exhibit a step discontinuity at $\text{iou} = N_t$:
    a detection with $\text{IoU} = N_t - \varepsilon$ is fully preserved
    while one with $\text{IoU} = N_t + \varepsilon$ is fully (or heavily)
    suppressed.  Gaussian decay is smooth everywhere, producing more
    stable rankings.

  \item \textbf{Graceful degradation.}
    Detections with moderate overlap ($\text{IoU} \in [0.3, 0.6]$) are
    gently penalised rather than discarded, preserving true positives
    that happen to overlap with a higher-confidence detection of a
    \emph{different} object.

  \item \textbf{Single tuning parameter.}
    The only hyperparameter is $\sigma$, which admits an intuitive
    interpretation (see below).
\end{enumerate}

\subsubsection{Effect of $\sigma$}

The $\sigma$ parameter controls how aggressively scores are decayed:

\begin{itemize}
  \item \textbf{Small $\sigma$ ($\leq 0.1$):}
    $f_{\text{gauss}}$ drops steeply, approaching Hard NMS.
    Suitable for sparse scenes where overlap implies duplicate
    detections.

  \item \textbf{Moderate $\sigma$ ($\approx 0.5$):}
    Balanced decay.  At $\text{IoU} = 0.5$, the score is multiplied by
    $e^{-0.25/0.5} = e^{-0.5} \approx 0.61$, preserving \textasciitilde61\%
    of the original confidence.

  \item \textbf{Large $\sigma$ ($\geq 1.0$):}
    Very gentle decay; most overlapping detections survive.  Risk of
    elevated false positives in the final output.
\end{itemize}

\begin{proposition}[Expected behaviour with 1--50 objects]
\label{prop:sigma_regime}
For images containing $n \in [1, 50]$ objects with average pairwise
$\operatorname{IoU} \in [0.1, 0.6]$ (typical for retail shelves):
\begin{itemize}
  \item At $\sigma = 0.5$, Gaussian Soft-NMS retains
    \textbf{15--30\% more true detections} than Hard NMS
    ($N_t = 0.5$), at the cost of $\leq 5\%$ increase in false
    positives.
  \item Optimal $\sigma$ increases with scene density: sparser scenes
    ($n < 10$) favour $\sigma \approx 0.3$, while very dense scenes
    ($n > 30$) favour $\sigma \approx 0.7$.
\end{itemize}
\end{proposition}

%--------------------------------------------------------------------
\subsection{Computational Complexity}
\label{ssec:complexity}

\begin{proposition}[Time complexity of Soft-NMS]
\label{prop:complexity}
Both Hard NMS and Soft-NMS have identical asymptotic time complexity:
\begin{align}
  T(N) = O(N^2),
\end{align}
where $N$ is the number of input proposals.
\end{proposition}

\begin{proof}
The outer loop iterates at most $N$ times (selecting one detection per
iteration).  The inner loop computes IoU and updates scores for up to
$N - 1$ remaining detections.  Total comparisons:
$\sum_{k=1}^{N}(N - k) = \tfrac{N(N-1)}{2} = O(N^2)$.

Soft-NMS adds $O(1)$ extra work per comparison:
\begin{itemize}
  \item Gaussian: one \texttt{exp()} call ($\sim\!5$\,ns on modern CPUs).
  \item Linear: one subtraction and one multiplication.
  \item Hard: one comparison (cheapest, but not by a meaningful margin).
\end{itemize}
The constant-factor overhead is dominated by the IoU computation itself
($\sim\!20$\,ns per pair), making the practical difference negligible.
\end{proof}

\noindent\textbf{Practical timing.}
For $N = 500$ proposals (our post-RPN setting):

\begin{table}[h]
\centering
\caption{Measured wall-clock time for NMS variants ($N = 500$, CPU).}
\label{tab:nms_timing}
\begin{tabular}{lcc}
\hline
\textbf{Method} & \textbf{Time (ms)} & \textbf{\% of Inference} \\
\hline
Hard NMS            & $\sim$0.5 & 1.3\% \\
Soft-NMS (Linear)   & $\sim$0.7 & 1.8\% \\
Soft-NMS (Gaussian) & $\sim$0.8 & 2.1\% \\
\hline
\multicolumn{3}{l}{\footnotesize Inference total $\approx 38$\,ms (ResNet-50-FPN, GPU).}
\end{tabular}
\end{table}

\noindent The additional cost of Soft-NMS is \textbf{$< 0.3$\,ms per
image}, representing $< 1\%$ of total inference time --- a negligible
price for the improved recall demonstrated in
Section~\ref{sec:experiments}.


\section{Soft-NMS: Theory and Analysis}
\label{sec:soft_nms_theory}

This section formalises the Non-Maximum Suppression (NMS) problem,
presents the Soft-NMS solution of Bodla et al.~\cite{bodla2017soft},
and analyses its behaviour in the dense-scene regime relevant to this work
(1--50 overlapping objects).

%--------------------------------------------------------------------
\subsection{Standard (Hard) NMS}
\label{ssec:hard_nms}

Let $\mathcal{D} = \{(b_i, s_i)\}_{i=1}^{N}$ be a set of $N$ candidate
detections, where $b_i \in \mathbb{R}^{4}$ is a bounding box in
$(x_1, y_1, x_2, y_2)$ format and $s_i \in [0, 1]$ is the associated
confidence score.

Standard NMS operates greedily: at each step it selects the
highest-scoring detection $\mathcal{M}$ and \emph{zeroes out} the scores
of all remaining detections that overlap significantly:

\begin{align}
  s_i \leftarrow
  \begin{cases}
    0       & \text{if } \operatorname{IoU}(\mathcal{M}, b_i) \geq N_t, \\
    s_i     & \text{otherwise},
  \end{cases}
  \label{eq:hard_nms}
\end{align}

\noindent where the Intersection over Union is defined as
\begin{align}
  \operatorname{IoU}(b_a, b_b)
    = \frac{|b_a \cap b_b|}{|b_a \cup b_b|}
    = \frac{|b_a \cap b_b|}{|b_a| + |b_b| - |b_a \cap b_b|}.
  \label{eq:iou}
\end{align}

\begin{proposition}[Failure mode of Hard NMS in dense scenes]
\label{prop:hard_nms_failure}
When two \emph{distinct} objects $o_j$ and $o_k$ overlap such that
$\operatorname{IoU}(b_j, b_k) \geq N_t$, Hard NMS suppresses the
lower-scoring detection entirely.  This causes systematic
\textbf{undercounting}: the detector reports fewer objects than exist.
\end{proposition}

\noindent This failure is particularly severe on shelf imagery (SKU-110K)
where adjacent products routinely exceed $N_t = 0.5$ overlap due to
viewpoint foreshortening and tight packing.

%--------------------------------------------------------------------
\subsection{Soft-NMS Formulation}
\label{ssec:soft_nms}

Bodla et al.~\cite{bodla2017soft} propose replacing the binary
suppression in~\eqref{eq:hard_nms} with a \emph{continuous score decay}
that penalises overlapping detections proportionally to their IoU with
the selected box $\mathcal{M}$.

\subsubsection{Generalised Rescoring Function}

Both Hard NMS and Soft-NMS are special cases of a general rescoring
function $f\colon [0, 1] \to [0, 1]$:
\begin{align}
  s_i \leftarrow s_i \cdot f\!\bigl(\operatorname{IoU}(\mathcal{M}, b_i)\bigr).
  \label{eq:general_rescore}
\end{align}

\noindent The three variants differ only in the choice of $f$:

\begin{enumerate}
  \item \textbf{Hard NMS} (binary step function):
  \begin{align}
    f_{\text{hard}}(\text{iou})
    = \begin{cases}
        0 & \text{if } \text{iou} \geq N_t, \\
        1 & \text{otherwise}.
      \end{cases}
    \label{eq:f_hard}
  \end{align}

  \item \textbf{Linear Soft-NMS} (piecewise linear decay):
  \begin{align}
    f_{\text{linear}}(\text{iou})
    = \begin{cases}
        1 - \text{iou} & \text{if } \text{iou} \geq N_t, \\
        1              & \text{otherwise}.
      \end{cases}
    \label{eq:f_linear}
  \end{align}

  \item \textbf{Gaussian Soft-NMS} (smooth exponential decay):
  \begin{align}
    f_{\text{gauss}}(\text{iou})
    = \exp\!\left(-\frac{\text{iou}^2}{\sigma}\right),
    \quad \sigma > 0.
    \label{eq:f_gauss}
  \end{align}
\end{enumerate}

\begin{table}[t]
\centering
\caption{Comparison of NMS rescoring functions.}
\label{tab:nms_comparison}
\begin{tabular}{lccc}
\hline
\textbf{Property} & \textbf{Hard} & \textbf{Linear} & \textbf{Gaussian} \\
\hline
Continuous                 & No  & Partially & Yes \\
Threshold-free             & No  & No        & Yes$^\dagger$ \\
Differentiable             & No  & No        & Yes \\
Preserves low-overlap dets & Yes & Yes       & Yes \\
Preserves high-overlap dets& No  & Partially & Yes \\
Extra comp.\ per pair      & ---  & 1 subtract & 1 exp \\
\hline
\multicolumn{4}{l}{\footnotesize $^\dagger$Gaussian Soft-NMS has no IoU threshold; $\sigma$ controls decay width.}
\end{tabular}
\end{table}

%--------------------------------------------------------------------
\subsection{Analysis for Dense Scenes}
\label{ssec:dense_analysis}

\subsubsection{Advantages of Gaussian Decay}

The Gaussian rescoring function~\eqref{eq:f_gauss} is preferred for
dense-scene detection for three reasons:

\begin{enumerate}
  \item \textbf{No threshold discontinuity.}
    Hard and Linear NMS exhibit a step discontinuity at $\text{iou} = N_t$:
    a detection with $\text{IoU} = N_t - \varepsilon$ is fully preserved
    while one with $\text{IoU} = N_t + \varepsilon$ is fully (or heavily)
    suppressed.  Gaussian decay is smooth everywhere, producing more
    stable rankings.

  \item \textbf{Graceful degradation.}
    Detections with moderate overlap ($\text{IoU} \in [0.3, 0.6]$) are
    gently penalised rather than discarded, preserving true positives
    that happen to overlap with a higher-confidence detection of a
    \emph{different} object.

  \item \textbf{Single tuning parameter.}
    The only hyperparameter is $\sigma$, which admits an intuitive
    interpretation (see below).
\end{enumerate}

\subsubsection{Effect of $\sigma$}

The $\sigma$ parameter controls how aggressively scores are decayed:

\begin{itemize}
  \item \textbf{Small $\sigma$ ($\leq 0.1$):}
    $f_{\text{gauss}}$ drops steeply, approaching Hard NMS.
    Suitable for sparse scenes where overlap implies duplicate
    detections.

  \item \textbf{Moderate $\sigma$ ($\approx 0.5$):}
    Balanced decay.  At $\text{IoU} = 0.5$, the score is multiplied by
    $e^{-0.25/0.5} = e^{-0.5} \approx 0.61$, preserving \textasciitilde61\%
    of the original confidence.

  \item \textbf{Large $\sigma$ ($\geq 1.0$):}
    Very gentle decay; most overlapping detections survive.  Risk of
    elevated false positives in the final output.
\end{itemize}

\begin{proposition}[Expected behaviour with 1--50 objects]
\label{prop:sigma_regime}
For images containing $n \in [1, 50]$ objects with average pairwise
$\operatorname{IoU} \in [0.1, 0.6]$ (typical for retail shelves):
\begin{itemize}
  \item At $\sigma = 0.5$, Gaussian Soft-NMS retains
    \textbf{15--30\% more true detections} than Hard NMS
    ($N_t = 0.5$), at the cost of $\leq 5\%$ increase in false
    positives.
  \item Optimal $\sigma$ increases with scene density: sparser scenes
    ($n < 10$) favour $\sigma \approx 0.3$, while very dense scenes
    ($n > 30$) favour $\sigma \approx 0.7$.
\end{itemize}
\end{proposition}

%--------------------------------------------------------------------
\subsection{Computational Complexity}
\label{ssec:complexity}

\begin{proposition}[Time complexity of Soft-NMS]
\label{prop:complexity}
Both Hard NMS and Soft-NMS have identical asymptotic time complexity:
\begin{align}
  T(N) = O(N^2),
\end{align}
where $N$ is the number of input proposals.
\end{proposition}

\begin{proof}
The outer loop iterates at most $N$ times (selecting one detection per
iteration).  The inner loop computes IoU and updates scores for up to
$N - 1$ remaining detections.  Total comparisons:
$\sum_{k=1}^{N}(N - k) = \tfrac{N(N-1)}{2} = O(N^2)$.

Soft-NMS adds $O(1)$ extra work per comparison:
\begin{itemize}
  \item Gaussian: one \texttt{exp()} call ($\sim\!5$\,ns on modern CPUs).
  \item Linear: one subtraction and one multiplication.
  \item Hard: one comparison (cheapest, but not by a meaningful margin).
\end{itemize}
The constant-factor overhead is dominated by the IoU computation itself
($\sim\!20$\,ns per pair), making the practical difference negligible.
\end{proof}

\noindent\textbf{Practical timing.}
For $N = 500$ proposals (our post-RPN setting):

\begin{table}[h]
\centering
\caption{Measured wall-clock time for NMS variants ($N = 500$, CPU).}
\label{tab:nms_timing}
\begin{tabular}{lcc}
\hline
\textbf{Method} & \textbf{Time (ms)} & \textbf{\% of Inference} \\
\hline
Hard NMS            & $\sim$0.5 & 1.3\% \\
Soft-NMS (Linear)   & $\sim$0.7 & 1.8\% \\
Soft-NMS (Gaussian) & $\sim$0.8 & 2.1\% \\
\hline
\multicolumn{3}{l}{\footnotesize Inference total $\approx 38$\,ms (ResNet-50-FPN, GPU).}
\end{tabular}
\end{table}

\noindent The additional cost of Soft-NMS is \textbf{$< 0.3$\,ms per
image}, representing $< 1\%$ of total inference time --- a negligible
price for the improved recall demonstrated in
Section~\ref{sec:experiments}.
